% ------------------------------------------------------------------------
%
\begin{frame}{Inductive proofs}

Let us remark that \(\fun{cat}([1], [2,3,4]) \twoheadrightarrow
[1,2,3,4] \twoheadleftarrow \fun{cat}([1,2],
[3,4])\).

\bigskip

It is enlightening to create \emph{equivalence classes} of terms that
are joinable. These classes then define an equivalence
relationship~\((\equiv)\) as follows:
\begin{center}
  \(a \equiv b\) if there exists a unique value~\(v\) such that \(a
  \xrightarrow{\smash[t]{*}} v\) and \(b \xrightarrow{\smash{*}} v\).
\end{center}
For instance, \(\fun{cat}([1,2], [3,4]) \equiv \fun{cat}([1],
[2,3,4])\).

\bigskip

The relation (\(\equiv\)) is indeed an equivalence because it is
\begin{itemize}

  \item \emph{reflexive}: \(a \equiv a\);

  \item \emph{symmetric}: if \(a \equiv b\), then \(b \equiv a\);

  \item \emph{transitive}: if \(a \equiv b\) and \(b \equiv c\), then
    \(a \equiv c\).

\end{itemize}

\end{frame}

% ------------------------------------------------------------------------
%
\begin{frame}{Inductive proofs}

  Of some interest are the following facts.
  \begin{itemize}

    \item If \(x_1 \leftarrow x_2 \twoheadrightarrow x_3 \equiv x_4
      \rightarrow x_5 \twoheadleftarrow x_6\), then \(x_1 \equiv x_2
      \equiv x_3 \equiv x_4 \equiv x_5 \equiv x_6\).

    \item If \(f(x)\) and~\(f(y)\) have a value, then
      \[x \equiv y \Rightarrow f(x) \equiv f(y),\]
      that is, \((\equiv)\) is a \textbf{congruence relation}, as a
      consequence of confluence and termination.

  \end{itemize}

  If we want to prove equivalences with variables ranging over
  \textbf{infinite sets}, like \(\fun{cat}(s,\fun{cat}(t,u)) \equiv
  \fun{cat}(\fun{cat}(s,t),u)\), we need the \textbf{induction
    principle}.

\end{frame}

% ------------------------------------------------------------------------
%
\begin{frame}{Well-founded induction}

We define a \textbf{well\hyp{}founded order} on a set~\(A\) as being a
binary relation \((\succ)\) such that it is
\begin{itemize}

  \item irreflexive: \(\forall x \in A. \neg(x \succ x)\),

  \item transitive: \(\forall x, y, z \in A. x \succ y \,\wedge\, y
    \succ z \Rightarrow x \succ z\),

  \item and does not have any \textbf{infinite descending chains}, to
    wit, there is no \(a_0 \succ a_1 \succ a_2 \succ \dots\)

\end{itemize}
The \textbf{well\hyp{}founded induction principle} then states that, for
any predicate~\(\aleph\) on a set~\(A\),
\begin{center}
  \(\forall a \in A.\aleph(a)\) is implied by \(\forall a.(\forall b.a
  \succ b \Rightarrow \aleph(b)) \Rightarrow \aleph(a)\).
\end{center}
When \(A=\mathbb{N}\), this principle is called \textbf{mathematical
(complete) induction}.

\end{frame}

% ------------------------------------------------------------------------
%
\begin{frame}{Well-founded induction}

\textbf{Structural induction} is another particular case where \(t \succ
s\) holds if, and only if, \(s\)~is a proper \textbf{subterm} of~\(t\),
namely, the abstract syntax tree of~\(s\) is included in the tree
of~\(t\) and \(s \neq t\).

\bigskip

Sometimes, a restricted form of structural induction is enough. For
instance, we can define \(\cons{x}{s} \succ s\), for any term~\(x\)
and any stack~\(s \in S\). Both \(x\)~and~\(s\) are \textbf{immediate
subterms} of \(\cons{x}{s}\). (In other words, we dispensed with the
transitivity axiom.) Because there are no infinite descending chains,
any subset \(B \subseteq A\) contains minimal elements \(M \subseteq
B\), that is, there is no \(b \in B\) such that \(a \succ b\), if \(a
\in M\). So
\begin{enumerate}

  \item we have to prove \(\aleph(a)\) for all \(a \in M\), the
    so\hyp{}called \textbf{basis}: here, only \(\aleph(\el)\), as
    \(\el\)~is by construction the unique, minimal stacks, and,

  \item assuming the so\hyp{}called \textbf{induction hypothesis}
    \(\aleph(s)\), we must prove \(\aleph(\cons{x}{s})\) for
    all~\(x\).

\end{enumerate}

\end{frame}
